\cotefile{83}
\scanpage{178}
\unsure{les homomorphismes voulus|l'homomorphisme voulu} de
\[
\deleted{\unsure{\mathop{*}_{i}|\bigotimes_i|*}}\, E_1\mathbin{\underset{\lambda}{*}}E_2\quad\text{dans}\quad E_1^\circ * E_2^\circ .
\]

Par passage à la limite etc., on en conclut l'injectivité si les $F_i$ sont plats.

Il faut définir un $\lambda$-scindage \inserted{\uline{canonique}} sur $* E_i$, donc il faut définir
\[
\mathop{*}^{i}E_i \longrightarrow \bigotimes E_i^\circ
\]
qui sur $\bigotimes E_i^\circ$ prenne la valeur $\lambda\,\mathrm{id}$. Il faut donc trouver un $\lambda$-accouplement can. à valeurs dans $\bigotimes E_i^\circ$.

On songe \deleted{tout} de suite à
\[
\begin{aligned}
f=\bigotimes u_i &: \bigotimes E_i\longrightarrow\bigotimes E_i^\circ\\
f(x_1,\ldots x_n)&=u_1(x_1)\otimes\ldots\otimes u_n(x_n)
\end{aligned}
\]
% slanted
\margin{mais il faudrait remplacer $\bigotimes E_i$ par son image -- ou bien ne pas \unsure{exiger|croire} dans les données $F_i\xrightarrow{\unsure{R_i|\beta_i}}E_i$ que celles-ci \unsure{sont|soient} injectives mais \unsure{seulement} donner un $u_i:E_i\mathbin{\unsure{-|\to}}F_i$ avec $u_i R_i=\lambda\,\mathrm{id}_{F_i}$.}
\[
\begin{aligned}
 f(x_1+\alpha_1,x_2\ldots x_n)&=(u_1(x_1)+\lambda\alpha_1)\otimes u_2(x_2)\ldots\otimes\unsure{u_i|u_n}(x_n)\\
 &=f(x_1,\ldots x_n)+\underset{\deleted{\unsure{g(x_2)}}}{\underline{\lambda\alpha_1}}\otimes u_2(x_2)\ldots u_n(x_n)
\end{aligned}
\]
prendre donc comme application associée
\[
 f^\circ(\alpha_1,\ldots,\alpha_n)=\lambda\alpha_1\otimes\ldots\otimes\alpha_n
\]
de sorte que les termes \unsure{en $\alpha$'s|non nuls} précédents \unsure{s'écrivent|s'écrive}
% archive sheet number: 177
